\mode*
\section{Search and verification protocol}
\label{app:protocol}

This appendix documents the literature searches behind the references, for
reproducibility, following the protocol of the companion paper
\autocite[App.~literature protocol]{VTProgMisconceptions}: every reference
was found by a recorded query, read at least at abstract level, and
recorded with a provenance block (claim, query, selection rationale,
verbatim supporting quote, verification level) directly above its entry in
\texttt{misconceptions.bib}.

\subsection{Phase one: seed round (2026-07-18)}

The seed round was a claim-driven search conducted with the
\texttt{scholar} tool against DBLP, OpenAlex, and Semantic Scholar
(Semantic Scholar was rate-limited during part of the session), with
abstracts fetched via the OpenAlex API and full texts from open repository
copies (the Isomöttönen \& Cochez author copy at the JYU repository was
read in full; the Gitless preprint is open at MIT DSpace).
The queries whose results were retained for this paper:

\begin{itemize}
  \item \texttt{scholar search "git misconceptions" -p dblp -p openalex}
  \item \texttt{scholar search "challenges confusions version control" -p dblp}
  \item \texttt{scholar search "purposes concepts misfits redesign git" -p dblp}
  \item \texttt{scholar search "shell tutor" -p dblp}
  \item \texttt{scholar search "black box inside glass box presenting computing concepts" -p openalex -p dblp}
  \item \texttt{scholar search "Sorva notional machines introductory programming" -p dblp -p openalex}
  \item \texttt{scholar search "Ben-Ari constructivism" -p dblp}
  \item \texttt{scholar search "difficulties of learning to program du Boulay" -p openalex}
  \item \texttt{scholar search '"variation theory" computer programming' -p openalex}
\end{itemize}

Discarded in this round: Feliciano et al.\ 2016 (student GitHub
experiences; abstract truncated, redundant next to the retained sources
--- to be re-evaluated in phase two); \enquote{Teaching Git on the side}
(ITiCSE 2015; platform paper, not misconceptions).
Not yet verified at full text (tracked in the course repository's issues
\#124--\#125): \autocite{DuBoulay1981BlackBox},
\autocite{BenAri1998Constructivism}.
Note also that \autocite{Postner2026GitMisconceptions} is a one-page
poster: low evidence weight, used as a pointer to the GitKit line of work.

\subsection{Phase two: systematic round (2026-07-18)}

The systematic round extended the seed search along the gaps it had left
open, using the same protocol: the \texttt{scholar} tool against DBLP and
OpenAlex (Semantic Scholar was rate-limited for part of the session),
abstracts fetched via the OpenAlex API, and full texts from open repository
copies where available.
Screening kept a result only if (i)~it concerned learning, teaching, or
documented use of Git or version control (not a platform- or tooling-only
paper), and (ii)~its abstract or full text was accessible so that a verbatim
supporting quote could be recorded; every retained reference carries a
provenance block in \texttt{misconceptions.bib}.

The productive queries, by gap:
\begin{description}
  \item[Version-control misconceptions] \texttt{scholar search "learning git
    version control difficulties" -p dblp -p openalex} retained
    \autocite{Yang2022GitCommands}, a large-scale Stack Overflow study of Git
    command difficulty.
  \item[Repository-data studies of student use] \texttt{scholar search
    "mining student version control repositories novices" -p dblp -p
    openalex}, with the title confirmed through a Crossref bibliographic
    query, retained the open-access \autocite{Cochez2013DVCSUsage}.
  \item[GitHub in education and the GitKit line] \texttt{scholar search
    "GitKit free open source software learning git" -p dblp -p openalex} and
    \texttt{scholar search "teaching git" -p dblp} retained the full GitKit
    paper \autocite{Braught2024GitKit}; the seed round had only the poster
    \autocite{Postner2026GitMisconceptions}.
  \item[Visualisation interventions] the same GitKit query surfaced, and we
    retained, \autocite{Bucher2026GitTakesTwo}, a controlled study of a
    split-view Git-learning tool.
  \item[Full-text re-read of the seed source] the open author copy of
    \autocite{IsomottonenCochez2014} was re-read in full; its
    Sections~5.1--5.7 supplied precise quotes for the staging-area
    (\texttt{add}-as-create-file), remote-prefix, Git/Bash-mixing, and
    blind-testing difficulties, which now back
    \cref{sec:model,sec:staging,sec:remotes} and are recorded as additional
    quotes in the entry's provenance block.
\end{description}

Re-evaluated and \emph{not} retained.
The seed round's discarded candidate---Feliciano, Storey \& Zagalsky's case
study of students' GitHub experiences (ICSE-C 2016, DOI
10.1145/2889160.2889195)---was re-examined as promised: its abstract is
available only in truncated form and the paper is not open access, so its
specific findings could not be verified; since the retained GitKit and
Isomöttönen--Cochez sources already cover GitHub in education, it is left
uncited rather than cited from an inaccessible abstract.
Also screened out: the Onward!\ essay \enquote{What's wrong with git?} (2013,
superseded by the retained Gitless paper \autocite{DeRossoJackson2016Gitless});
the GitKit tutorial and demo abstracts (2022 and 2025, redundant with
\autocite{Braught2024GitKit}); the IASGE \enquote{behavioral approach} survey
(HICSS 2021), whose scope is scholarly researchers and repository
preservation rather than course misconceptions; and several Git-metrics
papers on measuring student contributions, which concern grading rather than
misconceptions.

Unproductive queries, recorded for reproducibility (no new retained
reference): \texttt{scholar search "teaching version control students"},
\texttt{"student git commit behavior repository mining"}, \texttt{"git
visualization tool teaching learning"}, \texttt{"visualizing commit graph
learning git"}, and \texttt{"pull request workflow teaching students"} (all
\texttt{-p dblp -p openalex}) returned only off-topic or education-general
results.
The research-question mode (\texttt{scholar rq}) could not be driven
non-interactively in this session---it requires a classification step that
returned no output---so the round relied on direct \texttt{scholar search}
queries; the one research-question-generated query that did surface,
\texttt{GitHub software engineering education students challenges} (DBLP), hit
a transient server error.

Remaining gaps.
Pull-request-workflow studies in teaching, and a controlled comparison of
visualisation interventions, were not found at citable quality; and the
course's own grading-data misconceptions (\cref{sec:model}, issue~\#2) are
still to be added.
These are carried as open items rather than filled from unverifiable sources.

\subsection{Extended-database round (2026-07-19)}

The rounds above searched only the keyless providers (DBLP, OpenAlex,
Semantic Scholar); Web of Science, Scopus and IEEE~Xplore were configured
but unused.
This round repeated the version-control queries across all providers:
\texttt{scholar search "git version control students misconceptions
learning" -p wos -p scopus -p ieee} and \texttt{"version control
misconceptions students" -p wos -p scopus} re-surfaced
\autocite{IsomottonenCochez2014} and teaching-platform papers but no
misconception study beyond those already retained --- evidence that the
retained set covers what the added databases index.
Natural-language queries rank noticeably worse in Web of Science and
Scopus than in the CS-specific indexes; a future round should use their
field-tagged syntax (\texttt{scholar syntax}).

\subsection{Named-session field-tagged round (2026-07-19)}

This round re-ran the version-control searches with field-tagged syntax
(WoS \texttt{TS=}, Scopus \texttt{TITLE-ABS-KEY}), recorded under the
\texttt{scholar} session \texttt{vt-git-systematic} and exported in full
(queries, screened papers, keep/discard decisions with reasons) to
\texttt{literature-review/} in this repository.
Retained: \autocite{Hamer2021GitMetrics} (repository-data measurement of
student contributions; the methodological precedent for our
assignments-as-data design) and \autocite{Ardimento2022ProcessMining}
(process mining of IDE and VCS interaction logs).
Additionally, the terminal paper's SSH query (session
\texttt{vt-terminal-systematic}) surfaced
\autocite{Wagner2025RethinkGit}, a teaching tip documenting that students
report being overwhelmed by Git's command surface and recommending
model-first teaching; it is retained here.
No further misconception study appeared: the misconception-specific
retained set remains that of the earlier rounds.

\subsection{Measurement round (2026-08-23): has learners' understanding of
Git been measured, and which aspects does the literature say are hard?}
\label{sec:protocol-measurement}

The rounds above asked what the literature says about learners'
misconceptions of Git; this round asks two prior questions that
\cref{sec:method-instruments,sec:related} and the aspect lists of
\cref{sec:model,sec:staging,sec:branches,sec:remotes} rely on.
First, \emph{has learners' understanding of version control with
Git---the commit model, the staging area, branches and merging, remotes
and collaboration---actually been measured}, by a phenomenographic or
interview study of conceptions, by a diagnostic instrument or concept
inventory, or by a knowledge test, and what do the existing empirical
evaluations of Git teaching measure instead?
Second, \emph{which difficulties per objective does that empirical
literature document}, so that the candidate critical aspects can be read
off documented difficulties rather than posited?
A methodological sub-question is whether the pre-/post-test design of
\cref{sec:method-instruments} has a precedent in learning-study practice
and in the diagnostic-instrument literature; this part re-uses the
companion paper's answer \autocite{VTTerminal}.

\paragraph{Method.}
Twenty topic queries were run with the \texttt{scholar} tool under the
session \texttt{vt-git-measurement}, each against OpenAlex, Scopus and
Web of Science, with IEEE Xplore added to the Git-specific ones and DBLP
to the shorter ones (DBLP requires every term to match and returned
nothing for queries of five or more terms, which is recorded in the
export); Semantic Scholar was unavailable throughout the round
(\texttt{scholar providers check}: API key rejected, HTTP~403), so its
coverage is missing.
Twenty results per provider per query were kept.
The queries, verbatim:
\begin{itemize}
  \item \emph{Measurement of Git understanding.}
    \texttt{"git version control students misconceptions pre-test
    post-test"};
    \texttt{"concept inventory version control git"};
    \texttt{"phenomenographic study students understanding version
    control"};
    \texttt{"students mental model git commit branch interview study"};
    \texttt{"git knowledge test assessment students learning outcomes"};
    \texttt{"students conceptions of version control interview study"}.
  \item \emph{What the existing work measures.}
    \texttt{"teaching git course evaluation survey confidence students"};
    \texttt{"git learning tool visualization evaluation students"};
    \texttt{"novice git errors command usage analysis study"}.
  \item \emph{Difficulties per objective.}
    \texttt{"students understanding git commit history snapshot"};
    \texttt{"git staging area index students confusion"};
    \texttt{"git branching merging students difficulties
    misconceptions"};
    \texttt{"git push pull remote students confusion GitHub
    collaboration"};
    \texttt{"merge conflicts novices study version control"}.
  \item \emph{Method precedents} (the companion paper's queries, re-run so
    that the shared references are found in this session too):
    \texttt{"two-tier diagnostic test misconceptions development"};
    \texttt{"concept inventory computer science development validation
    interviews"};
    \texttt{"learning study pre-test post-test critical aspects Marton
    Pang"}.
  \item \emph{Counter-searches} (phrased to surface the interview and
    think-aloud studies that would refute the gap):
    \texttt{"ways of experiencing version control students conceptions
    interviews"};
    \texttt{"git mental models developers qualitative study"};
    \texttt{"version control think-aloud novices"}.
\end{itemize}
Because natural-language queries rank poorly in Web of Science and
Scopus (the lesson of the extended-database round above), two
field-tagged twins were run in each:
\texttt{TS=((git OR "version control") AND (student* OR novice* OR
learner*) AND (misconception* OR "mental model*" OR conception* OR
understanding) AND (test OR inventory OR interview* OR phenomenograph* OR
"pre-test" OR pretest OR "post-test"))} and
\texttt{TS=((git OR github OR "version control") AND (student* OR
novice* OR learner*) AND (difficult* OR misconception* OR confus* OR
barrier* OR error*) AND (commit* OR branch* OR merg* OR staging OR
remote* OR push OR pull OR clone))} in Web of Science, and their
\texttt{TITLE-ABS-KEY} counterparts in Scopus; and three short DBLP
queries, \texttt{"git misconceptions"}, \texttt{"git students"},
\texttt{"version control novices"}, since DBLP's implicit AND defeats the
long ones.
After the first hit lists had been read, eight complementary queries
widened the net (\texttt{"learning git novices"}; \texttt{"version
control education students understanding"}; \texttt{'"version control"
AND (misconception* OR "mental model*")'}; \texttt{"git mental model"};
\texttt{"novice programmers version control systems difficulties"};
\texttt{"git tutorial evaluation learning gains"}; \texttt{"hidden
dependencies version control user experience"}; \texttt{"students commit
messages behaviour analysis course"}, against OpenAlex with DBLP, IEEE
Xplore, Scopus or Web of Science as appropriate).
Their first run was lost from the session when a concurrently running
\texttt{scholar enrich} wrote the session file back---a tooling hazard
worth recording: \texttt{scholar} commands on one session must not run
concurrently---and was re-run; by then OpenAlex was answering
HTTP~429 (rate limit), so the re-run holds only the DBLP, IEEE, Scopus
and Web of Science hits, while the OpenAlex hits of the first run
survive in the saved CSV output and were screened from it by title (one
retained reference, \textcite{Church2014HiddenDependencies}, came from
that query and was recovered into the session through DBLP).
Seven references were not returned by the topic queries and were
retrieved as known items by title search (Scopus \texttt{TITLE()} or
DBLP), which their provenance blocks state: the concept-inventory review
and the operating-systems concept inventory
\autocite{Taylor2014CI,WebbTaylor2014OSCI} and the virtual-memory study
\autocite{Pamplona2024VirtualMemory}, shared with the companion paper;
and three Git-education papers identified through a web search for a
Git concept inventory---which found none, but surfaced a survey of Git as
a learning objective in statistics courses
\autocite{Beckman2021GitLearningObjective}, a pre/post survey evaluation
of a Git training programme \autocite{Jabrayilzade2022InteractiveGit},
and an author copy of the GitHub case study the seed round had left
uncited for want of access \autocite{Feliciano2016GitHubStudents}.
The session collected 725 distinct papers.
Abstracts were filled in with \texttt{scholar enrich} (259 added; 598 of
the 725 have one).
The 678 papers present at that point were classified zero-shot against a
written research context with \texttt{scholar llm classify} (GPT-5.4 via
GitHub Copilot), which kept 94 and discarded 584; the model's pass proved
a poor first sort for this question---of its 94 keeps only 10 survived
reading (it kept DevOps courses, think-aloud methodology and Copilot
studies), it discarded \textcite{IsomottonenCochez2014},
\textcite{DeRossoJackson2016Gitless} and the learning-study papers, and
its reason tags were misaligned with the papers on a number of records.
Every paper's decision and reason was therefore set by hand: the
Git-, version-control-, computing-education- and method-related titles
(about two hundred) were read with their abstracts, and the remainder
were assigned a coarse reason by reading their titles and venues; the
model's tags and confidences remain in each export record, preceding
our reason tag (the last tag of a record is ours), and its verdicts are
kept alongside ours in
\texttt{literature-review/vt-git-measurement-llm.csv} (reconstructed
from the model's kept list by title after our decisions had overwritten
its status field, so duplicate records of a kept title count as kept
there: 97 rather than 94).
The final tally is 31 kept (16 on Git, of which 7 were already cited
from the earlier rounds and are recorded as corroboration; 15 method
precedents, 7 of them cited) and 694 discarded, each with its reason.
Decisions were recorded with \texttt{scholar sessions decide}; the export
(queries, every screened paper with abstract, decision and reason) is in
\texttt{literature-review/vt-git-measurement.\{bib,csv,tex\}}, and each
retained reference carries a provenance block in
\texttt{misconceptions.bib} or \texttt{theory.bib}; the entries shared
with the companion paper carry their original blocks, with this paper's
retrieval and re-verification added.

\paragraph{Results.}
The hit list is dominated by false positives---science-education
diagnostic tests, nursing and psychology phenomenographies, software
repositories mined for other ends, and the many papers in which
\emph{commit}, \emph{branch} or \emph{merge} mean something else---and
the on-topic residue falls into five groups.

\emph{Conceptions studies.}
No phenomenographic or interview study of how learners understand Git
or version control exists; the counter-searches returned none, and the
phenomenographic work they did surface concerns nursing, mobile learning
and vectors.
The closest are three qualitative studies of a different kind:
interviews and observation of \emph{professional} engineers, who can
articulate Git's conceptual model yet follow narrow ritual workflows out
of fear, which the authors trace to Git's \emph{hidden
dependencies}---between the files on disk and the current branch, between
a local branch and the remote-tracking copy of the remote (a \texttt{git
status} reporting \enquote{up-to-date with origin/master} can be
hundreds of commits behind the remote), and between new files and the
index \autocite{Church2014HiddenDependencies}; interviews with students
about GitHub in two courses, in which students name unfamiliarity with
Git and GitHub as a challenge and ask for version control to be taught
earlier \autocite{Feliciano2016GitHubStudents}; and the teaching
observations and survey already cited \autocite{IsomottonenCochez2014}.
None collects students' own accounts of Git's objects.

\emph{Diagnostic instruments and concept inventories.}
No validated diagnostic instrument or concept inventory for Git or
version control was found, by the topic queries, by the field-tagged
queries or by a web search; the computing concept-inventory review
\autocite{Taylor2014CI} and the CS1 test items
\autocite{Vahrenhold2014TestItems} place the field's groundwork in
CS1 and digital logic, and the nearest instrument is the
operating-systems concept inventory \autocite{WebbTaylor2014OSCI}, as
for the shell.
What does exist is assessment practice: where Git is an explicit
learning objective, it is assessed through commit histories and rubric
points, and through in-class exam items of exactly the diagnostic shape
this paper uses---whether a pull modifies files on the local computer,
on the remote, both or neither \autocite{Beckman2021GitLearningObjective}.
The nearest \emph{measurement of conceptions} is again the
operating-systems study that reads conceptions of virtual memory off
written tests at the start and end of a course
\autocite{Pamplona2024VirtualMemory}.

\emph{Documented difficulties per objective}---the feed for the aspect
lists.
For the commit model: the commit-as-save conflation read off repository
data \autocite{Cochez2013DVCSUsage}, the blind-testing effect of a
missing mental model \autocite{IsomottonenCochez2014}, and the
\enquote{overwhelmed by commands} report \autocite{Wagner2025RethinkGit}.

For the staging area: \texttt{add} read as \emph{make} and the
\texttt{add .}/\texttt{commit -a} habit \autocite{IsomottonenCochez2014},
the index as a hidden dependency in which an unadded file is unknown to
Git \autocite{Church2014HiddenDependencies}, and the design-level misfit
analysis \autocite{DeRossoJackson2016Gitless}.

For branches: the branch-as-directory conception
\autocite{IsomottonenCochez2014} and the files-on-disk-depend-on-the-branch
dependency \autocite{Church2014HiddenDependencies}.

For remotes: the remote-prefix confusion \autocite{IsomottonenCochez2014},
the local-branch/remote-tracking/remote dependency and the misleading
\enquote{up-to-date} \autocite{Church2014HiddenDependencies}, the
distributed-state awareness a split-view tool sets out to build
\autocite{Bucher2026GitTakesTwo}, and the workflow misconceptions of the
GitKit line \autocite{Postner2026GitMisconceptions,Braught2024GitKit}.

A teacher focus group has produced a list of students' Git errors and
poor practices that is wider than any of these
\autocite{Eraslan2020GitErrors}; its full text is behind a paywall and
is cited at abstract scope.
% TODO(#5): read the Eraslan et al. 2020 full text through the library and
% add its per-objective items to the catalogue and the aspect lists.

\emph{What the evaluations measure.}
The rest of the empirical work measures perception, confidence,
performance or behaviour: pre- and post-surveys of self-assessed
sufficiency around a Git training programme (67\,\% to 9\,\%
\enquote{insufficient}) \autocite{Jabrayilzade2022InteractiveGit}, of
perspectives and engagement in a 319-student GitHub course
\autocite{Patani2024GitHubImpact}, of perceptions of Git concepts around
a tutorial \autocite{Sproull2022GitIt}, and a 315-student survey of how
students use GitHub \autocite{Tiwari2022GitHubSurvey}.

Skill and behaviour are the other two measures: a live skill checkoff
in which students demonstrate Git to a teaching assistant, with
confidence surveyed alongside \autocite{McMahon2026Checkoffs}; the
split-view tool's mixed performance gains
\autocite{Bucher2026GitTakesTwo}; and repository metrics and process
mining, for grading or for behaviour
\autocite{Hamer2021GitMetrics,Ardimento2022ProcessMining,Cochez2013DVCSUsage}---the
many further metrics-for-grading and GitHub-metrics-predict-performance
papers in the hit list were discarded as measuring contribution, not
understanding.

\emph{Learning-study precedents.}
As in the companion paper: the pre-test (written task plus interviews)
used to identify the critical aspects that separate the qualitatively
different ways of understanding, the same task repeated as post-test
\autocite{PangMarton2005}; the procedure as standard learning-study
practice \autocite{Kullberg2017ZDM}; and interviews and open-ended tests
as the most common route to students' conceptions in the
diagnostic-instrument literature
\autocite{KaltakciGurel2015Diagnostics}.

\paragraph{Conclusion.}
No: learners' understanding of Git has not been measured with an
instrument built for it, nor studied phenomenographically.
No concept inventory or validated diagnostic for Git or version control
exists, and no study collects students' own accounts of commits,
branches or remotes; what the empirical work measures is self-assessed
confidence, perception, skill performance and repository behaviour.
The difficulties, however, are better documented than for the shell:
course observation, a teacher focus group, repository data and a
usability analysis of professionals together name the commit-as-save
conflation, \texttt{add} as \emph{make} and the hidden index, the branch
as a directory and the files' dependency on the branch, and the
local/remote-tracking/remote confusion---so several of the candidate
aspects of \cref{sec:model,sec:staging,sec:branches,sec:remotes} can be
read off documented difficulties, which the aspect lists record, while
none rests on learners' own accounts.
The qualifications are that diagnostic-shaped exam items for Git exist
in teaching practice \autocite{Beckman2021GitLearningObjective}, that
conceptions have been read off written pre- and post-tests in the
neighbouring operating-systems course \autocite{Pamplona2024VirtualMemory},
and that the learning-study tradition
\autocite{PangMarton2005,Kullberg2017ZDM} and the diagnostic-instrument
typology \autocite{KaltakciGurel2015Diagnostics} give the pre-/post-test
design of \cref{sec:method-instruments} its precedent.
Semantic Scholar's absence, OpenAlex's rate limit on the re-run of the
complementary queries, and the paywalled teacher-focus-group list are the
round's limitations; the claim rests on OpenAlex, Scopus, Web of
Science, IEEE Xplore and DBLP.
